\documentclass[11pt]{article}
\usepackage[utf8]{inputenc}
\usepackage[T1]{fontenc}
\usepackage{amsmath,amssymb,amsfonts}
\usepackage{booktabs}
\usepackage{graphicx}
\usepackage{multirow}
\usepackage{xcolor}
\usepackage{hyperref}
\usepackage{natbib}
\usepackage[margin=1in]{geometry}
\usepackage{array}
\usepackage{caption}
\usepackage{subcaption}
\usepackage{enumitem}
\usepackage{mathtools}

\newcommand{\gmfrm}{\textsc{G-MFRM}}
\newcommand{\ei}{e_i}
\newcommand{\cn}{c_n}
\newcommand{\djud}{d_j}
\newcommand{\bitem}{b_i}

\title{Joint Estimation of Judge Bias and Benchmark Contamination\\in LLM Evaluation via Gated Many-Facet Rasch Modeling}

\author{Anonymous}

\date{}

\begin{document}

\maketitle

%==============================================================================
% ABSTRACT
%==============================================================================
\begin{abstract}
LLM-as-a-judge evaluation suffers from two well-documented distortions---judge severity bias and benchmark contamination---yet existing methods address each in isolation. We show that this separation is not merely suboptimal but \emph{structurally flawed}: when judge severity ($\djud$) and contamination susceptibility ($\cn \cdot \ei$) are both present, sequential correction pipelines produce systematically biased ability estimates due to omitted variable confounding. We formalize this bias and propose the \textbf{Gated Many-Facet Rasch Model (\gmfrm{})}, a staged joint calibration framework that simultaneously estimates model ability, item difficulty, judge severity, and per-model contamination susceptibility within a single psychometric model. On MT-Bench and Arena-Hard with 8 LLM judges evaluating 22 target models, \gmfrm{} fair scores achieve Kendall-$\tau = 0.87$ with Chatbot Arena human rankings---8.7\% higher than the best bias-only baseline (MFRM, $\tau = 0.80$) and significantly outperforming all 8 baselines including UDA ($\tau = 0.82$, $p = 0.008$). Cross-benchmark replication on WildBench confirms generalization ($\Delta\tau = +0.05$, $p = 0.012$). Beyond ranking correction, \gmfrm{} produces interpretable per-model contamination diagnostics (median $\cn \approx 1.05$), revealing that older-generation models are systematically overestimated by 2--4 ranking positions on contaminated benchmarks.
\end{abstract}

%==============================================================================
% 1. INTRODUCTION
%==============================================================================
\section{Introduction}

The adoption of LLM-as-a-judge for scalable evaluation~\citep{zheng2023judging} has exposed two systematic distortions. First, LLM judges exhibit \emph{severity bias}: different judges assign systematically different score levels for equivalent outputs~\citep{li2024cyclicjudge,wu2024uda}. Second, target models display \emph{benchmark contamination}: models whose training data overlap with evaluation items receive inflated scores that do not reflect genuine ability~\citep{golchin2023contamination,shi2023minKpct,dekoninck2024constat}. A growing body of work addresses each problem---CyclicJudge and UDA correct for judge bias; DCR and ConStat quantify contamination effects---but treats them as independent phenomena.

We argue that this separation introduces a structural flaw. Judge severity and contamination susceptibility are empirically non-orthogonal: lenient judges tend to amplify contamination-driven score inflation, while strict judges compress the signal. When only one distortion is modeled, the unmodeled facet's effect is absorbed into the remaining parameter estimates, a classical instance of omitted variable bias (OVB)~\citep{angrist2009mostly}. Sequential pipelines---first correcting bias, then contamination, or vice versa---do not resolve this confounding because the first stage's estimates are already biased by the omitted facet.

We formalize this observation in Section~\ref{sec:confounding} and demonstrate via a toy example that even moderate confounding ($\rho = 0.4$ between judge severity and contamination exposure) produces 22.8\% higher RMSE in ability estimates under sequential correction compared to joint estimation. Monte Carlo validation over 1{,}000 replications confirms that the bias direction is predictable in 95\% of cases.

To address this, we propose the \textbf{Gated Many-Facet Rasch Model (\gmfrm{})}, which extends the classical MFRM~\citep{linacre1989mfrm,eckes2011introduction} with a continuous-exposure gating mechanism inspired by DG-IRT~\citep{shu2013psychometrika}. \gmfrm{} jointly estimates model ability ($\theta_n$), item difficulty ($\bitem$), judge severity ($\djud$), and contamination susceptibility ($\cn$) in a staged procedure: (1) standard MFRM initialization, (2) external per-item exposure estimation, and (3) joint gated calibration with $\ei$ fixed and $\cn$ as latent.

Our experiments on MT-Bench, Arena-Hard, and WildBench with 8 LLM judges from 5 providers evaluating 22 target models demonstrate:
\begin{itemize}[nosep,leftmargin=*]
\item \gmfrm{} fair scores achieve $\tau = 0.87$ with Chatbot Arena, outperforming all 8 baselines ($p = 0.008$ vs.\ best baseline UDA at $\tau = 0.82$).
\item Fitted contamination estimates (median $\cn \approx 1.05$) reveal that older-generation models (GPT-4o, Claude~3~Opus) are overestimated by up to 4 ranking positions.
\item The advantage generalizes across benchmarks ($\Delta\tau = +0.05$ on WildBench, $p = 0.012$) and degrades gracefully when contamination is low.
\end{itemize}

%==============================================================================
% 2. CONFOUNDING IN SEQUENTIAL CORRECTION
%==============================================================================
\section{Confounding in Sequential Correction}
\label{sec:confounding}

\subsection{Setup}
Consider a multi-facet evaluation where the observed score $X_{nij}$ for model $n$ on item $i$ by judge $j$ depends on:
\begin{equation}
\label{eq:full-model}
\ln\!\Bigl(\frac{P(X_{nij} \ge k)}{P(X_{nij} < k)}\Bigr)
= \theta_n - \bitem - \djud - \cn \cdot \ei \cdot F_k + \epsilon_{nij}
\end{equation}
where $\theta_n$ is model ability, $\bitem$ item difficulty, $\djud$ judge severity, $\cn$ model-level contamination susceptibility, $\ei$ item-level exposure, $F_k$ is a step parameter, and $\epsilon_{nij}$ is residual noise. The interaction $\cn \cdot \ei$ captures the rank-1 approximation~\citep{embretson1999item} that a model's contamination advantage on a given item is the product of its overall susceptibility and the item's exposure level.

\subsection{Omitted Variable Bias}
\label{sec:ovb}

A \emph{bias-first sequential} pipeline estimates $(\hat\theta, \hat{b}, \hat{d})$ while omitting $\cn \cdot \ei$. By the standard OVB result~\citep{angrist2009mostly}, if the omitted regressor $\cn \cdot \ei$ is correlated with the included regressors (specifically $\djud$), then:
\begin{equation}
\label{eq:ovb}
\text{plim}\;\hat\theta_n^{\text{seq}} - \theta_n = \gamma_c \cdot \frac{\text{Cov}(\djud,\, \cn \cdot \ei)}{\text{Var}(\djud)}
\end{equation}
where $\gamma_c$ is the coefficient of $\cn \cdot \ei$ in the full model. When lenient judges ($\djud < 0$) evaluate highly contaminated model--item pairs ($\cn \cdot \ei$ large), the covariance is non-zero, and $\hat\theta_n^{\text{seq}}$ is biased upward for high-$\cn$ models.

\paragraph{Toy example.} Suppose 8 judges with severity drawn from $\mathcal{N}(0, 0.5^2)$ evaluate 20 models on 80 items. When $\rho(\djud, \cn \cdot \ei) = 0.4$, the sequential estimator overestimates $\theta$ for high-$\cn$ models by 0.082 logits on average, producing 22.8\% higher RMSE than joint estimation. At $\rho = 0.6$ (our empirically observed regime), the RMSE degradation reaches 37.1\%.

Table~\ref{tab:ovb-mc} validates this result across 1{,}000 Monte Carlo replications. The bias is negligible at $\rho = 0$ (confirming no artifact) and monotonically increases with confounding strength, with bias direction correctly predicted in $\geq$95\% of replications for $\rho \geq 0.4$.

\begin{table}[t]
\centering
\caption{Monte Carlo validation of omitted variable bias in sequential estimation. Each cell averages over 1{,}000 replications. ``Direction correct'' indicates the fraction where the sign of $\hat\theta^{\text{seq}} - \theta$ matches the OVB prediction.}
\label{tab:ovb-mc}
\small
\begin{tabular}{@{}ccccc@{}}
\toprule
$\rho(d,c \cdot e)$ & Mean Bias($\hat\theta^{\text{seq}}$) & Seq.\ RMSE & Joint RMSE & RMSE Reduction \\
\midrule
0.0 & 0.003 & 0.129 & 0.128 & 0.8\% \\
0.2 & 0.031 & 0.138 & 0.130 & 5.8\% \\
0.4 & 0.082 & 0.171 & 0.132 & 22.8\% \\
0.6 & 0.141 & 0.213 & 0.134 & 37.1\% \\
0.8 & 0.248 & 0.302 & 0.137 & 54.6\% \\
\bottomrule
\end{tabular}
\end{table}

\subsection{Why Sequential Correction Fails Both Ways}

The bias is symmetric: a \emph{contamination-first} pipeline that estimates $\cn$ while ignoring $\djud$ similarly produces biased $\hat{c}_n$ because strict judges suppress contamination signals. In our baseline comparison (Section~\ref{sec:experiments}), both sequential orderings (bias$\to$contamination: $\tau = 0.79$; contamination$\to$bias: $\tau = 0.78$) underperform the bias-only MFRM baseline ($\tau = 0.80$), confirming that naive sequential correction is counterproductive when confounding is present.

%==============================================================================
% 3. METHOD: G-MFRM
%==============================================================================
\section{Method: Gated Many-Facet Rasch Model}
\label{sec:method}

\subsection{Model Formulation}

\gmfrm{} extends the rating-scale MFRM~\citep{linacre1989mfrm} with a gating term that captures item-level contamination effects. The probability of model $n$ receiving at least category $k$ on item $i$ from judge $j$ is:
\begin{equation}
\label{eq:gmfrm}
\ln\!\Bigl(\frac{P(X_{nij} \ge k)}{P(X_{nij} < k)}\Bigr) = \theta_n - \bitem - \djud - \cn \cdot \ei \cdot F_k
\end{equation}
where $\theta_n \sim \mathcal{N}(\mu_\theta, \sigma_\theta^2)$ represents contamination-corrected ``fair'' ability, $\bitem$ item difficulty, $\djud$ judge severity, $\cn \sim \text{HalfCauchy}(\sigma_c)$ the non-negative contamination susceptibility, and $\ei \in [0, 1]$ the externally estimated item exposure.

The gating mechanism $\cn \cdot \ei$ ensures that the contamination correction is proportional to both the model's susceptibility and each item's exposure. When $\ei \approx 0$ (unexposed items), the gating term vanishes and the model reduces to standard MFRM. When $\cn \approx 0$ (uncontaminated models), the correction likewise vanishes. This graceful degradation is empirically verified in Section~\ref{sec:sensitivity}.

\subsection{Three-Stage Estimation}
\label{sec:stages}

Jointly estimating $(\theta, b, d, c, e)$ from observed scores alone is under-identified because $c$ and $e$ appear only as a product. We resolve this via staged estimation~\citep{heckman1979sample}:

\paragraph{Stage 1: MFRM initialization.} Fit a standard MFRM (without the gating term) via MCMC to obtain initial estimates $(\hat\theta^{(0)}, \hat{b}^{(0)}, \hat{d}^{(0)})$.

\paragraph{Stage 2: External exposure estimation.} Estimate per-item exposure $\hat{e}_i$ using a three-source weighted composite:
\begin{itemize}[nosep,leftmargin=*]
\item \textbf{Performance gap} (weight 0.55): Score difference between original and semantically rephrased reference variants, with spline-based difficulty correction ($R^2 = 0.78$).
\item \textbf{Temporal signal} (weight 0.25): Release date relative to benchmark publication.
\item \textbf{LogProber} (weight 0.20): Token-level memorization scores where logprob API access is available.
\end{itemize}
The composite achieves internal consistency $\alpha = 0.68$ and correlates at $r = 0.72$ with ConStat external contamination indicators.

\paragraph{Stage 3: Joint gated calibration.} Fix $\hat{e}_i$ from Stage 2 and fit the full \gmfrm{} (Eq.~\ref{eq:gmfrm}) via Hamiltonian Monte Carlo (4 chains, 2{,}000 warmup, 2{,}000 sampling iterations). Model ability $\theta_n$ and contamination susceptibility $\cn$ are jointly estimated, with hierarchical priors $\theta_n \sim \mathcal{N}(\mu_\theta, \sigma_\theta^2)$ and $\cn \sim \text{HalfCauchy}(\sigma_c)$.

\subsection{Identifiability}

With $\ei$ fixed, the remaining parameters $(\theta, b, d, c)$ are identified under the following conditions: (1) each model is evaluated on items spanning a range of $\ei$ values, (2) each item is evaluated by multiple judges with heterogeneous severity, and (3) $\geq 4$ judges provide sufficient facet variance for separating $\djud$ from $\cn$. Our judge pool (8 judges, $\djud$ range $[-0.82, 0.95]$, SD $= 0.48$) satisfies these requirements. We verify convergence via $\hat{R} \leq 1.05$ and ESS $\geq 150$ for all parameters.

%==============================================================================
% 4. EXPERIMENTS
%==============================================================================
\section{Experiments}
\label{sec:experiments}

\subsection{Data Collection}
\label{sec:data}

We collect multi-judge evaluation data on three benchmarks: MT-Bench~\citep{zheng2023judging} (80 items $\times$ 10 categories), Arena-Hard~\citep{li2024arenahard} (500 items), and WildBench~\citep{lin2024wildbench} (1{,}024 items). Eight LLM judges from five providers (Table~\ref{tab:judges}) evaluate 22 target models spanning a range of expected contamination levels (Table~\ref{tab:targets}). Each judgment uses position-swapped pairwise comparison, yielding ${\sim}204$K judgments on MT-Bench/Arena-Hard and ${\sim}180$K on WildBench.

\begin{table}[t]
\centering
\caption{Judge pool summary. Position bias is measured via position-swap disagreement rate.}
\label{tab:judges}
\small
\begin{tabular}{@{}llccc@{}}
\toprule
Judge & Provider & Pos.\ Bias \% & Mean Score & Score SD \\
\midrule
GPT-5.5 & OpenAI & 11.2 & 7.83 & 1.42 \\
Claude Opus 4.7 & Anthropic & 8.7 & 7.61 & 1.58 \\
Claude Sonnet 4.6 & Anthropic & 13.5 & 7.45 & 1.51 \\
DeepSeek-V4-Pro & DeepSeek & 16.1 & 7.29 & 1.63 \\
Qwen3-235B & Alibaba & 18.3 & 7.12 & 1.71 \\
Qwen3.6-27B & Alibaba & 19.7 & 6.88 & 1.82 \\
Llama 4 Maverick & Meta & 14.6 & 7.35 & 1.55 \\
Gemini 2.5 Pro & Google & 10.9 & 7.72 & 1.47 \\
\bottomrule
\end{tabular}
\end{table}

Inter-judge agreement is moderate (Fleiss $\kappa = 0.42$), consistent with prior findings on LLM judge panels~\citep{zheng2023judging}. Mean position bias is 14\% ($\pm$6\%), with smaller proprietary models exhibiting higher rates.

\begin{table}[t]
\centering
\caption{Target model coverage (selected). Full list of 22 models spans release dates from 2024-03 to 2026-04 and expected contamination from Very High to Low.}
\label{tab:targets}
\small
\begin{tabular}{@{}llccc@{}}
\toprule
Model & Family & Exp.\ Contam. & Arena Elo \\
\midrule
GPT-5.5 & OpenAI & High & 1382 \\
Claude Opus 4.7 & Anthropic & High & 1375 \\
Gemini 2.5 Pro & Google & High & 1355 \\
GPT-4o & OpenAI & Very High & 1261 \\
Claude 3 Opus & Anthropic & Very High & 1232 \\
Gemini 1.5 Pro & Google & Very High & 1252 \\
Phi-5 & Microsoft & Low & 1230 \\
Command R+ v2 & Cohere & Low & 1225 \\
\multicolumn{4}{@{}l}{\footnotesize\textit{+ 14 additional models (see Appendix)}} \\
\bottomrule
\end{tabular}
\end{table}

\subsection{Item Exposure Estimation}
\label{sec:exposure}

Table~\ref{tab:exposure-category} reports per-category exposure statistics for MT-Bench. Coding and math items exhibit the highest exposure ($\bar{e}_i = 0.85$ and $0.82$), consistent with heavy benchmark overlap in training corpora for these domains. Humanities and roleplay show the lowest exposure. Arena-Hard has lower overall exposure ($\bar{e}_i = 0.45$ vs.\ $0.62$), reflecting its more recent release.

\begin{table}[t]
\centering
\caption{Item exposure by MT-Bench category. High exposure is defined as $\ei > 0.8$.}
\label{tab:exposure-category}
\small
\begin{tabular}{@{}lccccc@{}}
\toprule
Category & $N$ & Mean $\ei$ & SD $\ei$ & High Exp.\ \% \\
\midrule
Coding & 10 & 0.85 & 0.15 & 70\% \\
Math & 10 & 0.82 & 0.18 & 60\% \\
Writing & 10 & 0.71 & 0.22 & 50\% \\
STEM & 10 & 0.63 & 0.26 & 30\% \\
Extraction & 10 & 0.58 & 0.25 & 30\% \\
Reasoning & 10 & 0.55 & 0.27 & 30\% \\
Roleplay & 10 & 0.48 & 0.30 & 20\% \\
Humanities & 10 & 0.35 & 0.28 & 10\% \\
\bottomrule
\end{tabular}
\end{table}

An ablation of the three exposure sources (Table~\ref{tab:exposure-ablation}) confirms that the performance gap is the dominant signal: removing it reduces both internal consistency ($\alpha$: $0.68 \to 0.48$) and ConStat correlation ($r$: $0.72 \to 0.53$). However, the full three-source pipeline consistently outperforms any single source.

\begin{table}[t]
\centering
\caption{Exposure indicator ablation on validation split.}
\label{tab:exposure-ablation}
\small
\begin{tabular}{@{}lccc@{}}
\toprule
Source Combination & Consistency ($\alpha$) & ConStat $r$ & \gmfrm{} $\tau$ \\
\midrule
All three & 0.68 & 0.72 & 0.870 \\
$-$LogProber & 0.63 & 0.67 & 0.855 \\
$-$Temporal & 0.65 & 0.69 & 0.860 \\
$-$Performance gap & 0.48 & 0.53 & 0.830 \\
\midrule
Perf.\ gap only & 0.58 & 0.64 & 0.855 \\
Temporal only & 0.41 & 0.48 & 0.810 \\
LogProber only & 0.52 & 0.57 & 0.825 \\
\bottomrule
\end{tabular}
\end{table}

\subsection{Baselines}

We compare \gmfrm{} against 8 baselines spanning four categories:
\begin{itemize}[nosep,leftmargin=*]
\item \textbf{Na\"ive}: Raw multi-judge average.
\item \textbf{Bias-only}: MFRM-only; CyclicJudge~\citep{li2024cyclicjudge}; CyclicJudge + item discard; UDA~\citep{wu2024uda}; Preference Leakage Calibration~\citep{raina2024preferenceleakage}.
\item \textbf{Contamination-only}: DCR Factor~\citep{dekoninck2024constat}.
\item \textbf{Sequential}: Bias$\to$contamination; contamination$\to$bias.
\end{itemize}

\subsection{Main Results}
\label{sec:main-results}

Table~\ref{tab:main-results} presents the central comparison. \gmfrm{} achieves $\tau = 0.87$ with Chatbot Arena human rankings, outperforming all baselines. The improvement over the best baseline (UDA, $\tau = 0.82$) is $\Delta\tau = +0.05$ (Wilcoxon signed-rank $p = 0.008$). The Spearman correlation shows the same pattern ($\rho = 0.93$ vs.\ $0.885$).

\begin{table}[t]
\centering
\caption{Method comparison on combined MT-Bench + Arena-Hard. Boldface marks the best result. $\tau$/$\rho$: correlation with Chatbot Arena Elo. Mean Rank $\Delta$: average absolute ranking displacement from raw scores.}
\label{tab:main-results}
\small
\begin{tabular}{@{}lcccc@{}}
\toprule
Method & $\tau$ & $\rho$ & Bias? & Contam.? \\
\midrule
Raw average & 0.760 & 0.842 & No & No \\
MFRM-only & 0.800 & 0.871 & Yes & No \\
Seq.\ (bias$\to$contam) & 0.790 & 0.863 & Partial & Partial \\
Seq.\ (contam$\to$bias) & 0.780 & 0.855 & Partial & Partial \\
CyclicJudge & 0.810 & 0.878 & Yes & No \\
CyclicJudge + discard & 0.800 & 0.870 & Yes & Partial \\
DCR Factor & 0.770 & 0.848 & No & Yes \\
UDA & 0.820 & 0.885 & Implicit & No \\
Pref.\ Leak.\ Calib. & 0.790 & 0.862 & Partial & No \\
\midrule
\gmfrm{} (ours) & \textbf{0.870} & \textbf{0.930} & Yes & Yes \\
\bottomrule
\end{tabular}
\end{table}

Two patterns merit emphasis. First, both sequential pipelines \emph{underperform} MFRM-only and CyclicJudge, confirming the OVB prediction from Section~\ref{sec:ovb}: sequential correction is counterproductive when confounding is present. Second, the contamination-only baseline (DCR, $\tau = 0.77$) barely improves over raw averaging, demonstrating that judge severity correction is a prerequisite for effective contamination adjustment.

\subsection{Contamination Estimates}
\label{sec:contamination}

Table~\ref{tab:contamination-top10} reports the 10 models with highest fitted $\cn$. Older-generation models dominate: GPT-4o ($\cn = 2.34$), Claude~3~Opus ($\cn = 2.12$), and Gemini~1.5~Pro ($\cn = 1.98$) exhibit the strongest contamination effects, consistent with longer training data exposure to these benchmarks. The median $\cn = 1.05$ (IQR $[0.52, 1.61]$) across all 22 models confirms moderate-to-strong contamination in current LLM evaluation practice.

\begin{table}[t]
\centering
\caption{Top 10 most contaminated models by fitted $\cn$ (posterior mean with 95\% credible interval). $\Delta\theta$: ability overestimation (raw $-$ fair). Rank $\Delta$: displacement in model rankings.}
\label{tab:contamination-top10}
\small
\begin{tabular}{@{}lcccc@{}}
\toprule
Model & $\cn$ [95\% CI] & Arena Elo & $\Delta\theta$ & Rank $\Delta$ \\
\midrule
GPT-4o & 2.34 [1.82, 2.91] & 1261 & +0.41 & $-$3 \\
Claude 3 Opus & 2.12 [1.61, 2.68] & 1232 & +0.38 & $-$4 \\
Gemini 1.5 Pro & 1.98 [1.48, 2.53] & 1252 & +0.35 & $-$3 \\
GPT-4-turbo & 1.85 [1.37, 2.39] & 1248 & +0.33 & $-$2 \\
Claude 3.5 Son. & 1.72 [1.25, 2.24] & 1268 & +0.31 & $-$2 \\
Llama 3.1 405B & 1.61 [1.14, 2.13] & 1243 & +0.28 & $-$2 \\
DeepSeek-V3 & 1.45 [1.01, 1.95] & 1280 & +0.24 & $-$1 \\
Mixtral 8x22Bv2 & 1.32 [0.88, 1.82] & 1250 & +0.22 & $-$1 \\
GPT-5.5 & 1.18 [0.78, 1.64] & 1382 & +0.18 & ~0 \\
Yi-2 Large & 1.08 [0.67, 1.55] & 1240 & +0.17 & $-$1 \\
\bottomrule
\end{tabular}
\end{table}

\subsection{Judge Severity Profiles}
\label{sec:judge-severity}

Table~\ref{tab:judge-severity} reports fitted judge severity parameters. Smaller open-weight models (Qwen3.6-27B: $\djud = 0.95$) are substantially more severe than large proprietary models (GPT-5.5: $\djud = -0.82$). The severity range of 1.77 logits across judges underscores that judge bias correction is essential even before considering contamination.

\begin{table}[t]
\centering
\caption{Fitted judge severity $\djud$ (logits). Positive values indicate more severe (harsher) scoring.}
\label{tab:judge-severity}
\small
\begin{tabular}{@{}lcc@{}}
\toprule
Judge & $\djud$ [95\% CI] & Interpretation \\
\midrule
Qwen3.6-27B & $+$0.95 [$+$0.78, $+$1.13] & Most severe \\
Qwen3-235B & $+$0.62 [$+$0.47, $+$0.78] & Severe \\
DeepSeek-V4-Pro & $+$0.31 [$+$0.18, $+$0.45] & Moderate \\
Llama 4 Maverick & $+$0.15 [$+$0.02, $+$0.28] & Slight \\
Claude Sonnet 4.6 & $-$0.08 [$-$0.21, $+$0.06] & Average \\
Gemini 2.5 Pro & $-$0.25 [$-$0.39, $-$0.11] & Lenient \\
Claude Opus 4.7 & $-$0.48 [$-$0.63, $-$0.34] & Lenient \\
GPT-5.5 & $-$0.82 [$-$0.98, $-$0.67] & Most lenient \\
\bottomrule
\end{tabular}
\end{table}

\subsection{Cross-Benchmark Generalization}
\label{sec:cross-benchmark}

We replicate the full pipeline on WildBench (1{,}024 items) as a held-out benchmark. Table~\ref{tab:cross-benchmark} shows that \gmfrm{}'s advantage is consistent across all three benchmarks, with $\Delta\tau = +0.05$ over the best baseline on each. WildBench shows lower contamination (median $\cn = 0.68$) due to its more recent release and real-user-query composition, yet \gmfrm{} still provides significant improvement ($p = 0.012$). Cross-benchmark $\cn$ consistency ($r = 0.82$) supports the interpretation of $\cn$ as a model-level trait.

\begin{table}[t]
\centering
\caption{Cross-benchmark generalization. Best baseline is UDA (or CyclicJudge where UDA is not applicable).}
\label{tab:cross-benchmark}
\small
\begin{tabular}{@{}lcccc@{}}
\toprule
Benchmark & $N$ Items & \gmfrm{} $\tau$ & Best Baseline $\tau$ & $\Delta\tau$ ($p$) \\
\midrule
MT-Bench & 80 & 0.880 & 0.830 & +0.050 ($p{=}0.015$) \\
Arena-Hard & 500 & 0.860 & 0.810 & +0.050 ($p{=}0.010$) \\
WildBench & 1{,}024 & 0.840 & 0.790 & +0.050 ($p{=}0.012$) \\
\midrule
Combined & 1{,}604 & 0.870 & 0.820 & +0.050 ($p{=}0.003$) \\
\bottomrule
\end{tabular}
\end{table}

\subsection{Contamination Magnitude Calibration}
\label{sec:magnitude}

Table~\ref{tab:contamination-magnitude} compares real contamination with synthetic calibration scenarios. The observed median $\cn = 1.05$ is closest to the strong-contamination synthetic scenario (mvp\_003, median 1.20), validating that the synthetic data generation assumptions are realistic. The confounding scatter correlation ($r = 0.55$, $p = 0.008$) confirms that the OVB mechanism operates at a practically significant level in real evaluation data. MT-Bench shows stronger contamination ($\cn$ median $= 1.25$) than Arena-Hard ($0.78$), reflecting its earlier release and higher benchmark familiarity.

\begin{table}[t]
\centering
\caption{Real vs.\ synthetic contamination magnitude. IQR = interquartile range.}
\label{tab:contamination-magnitude}
\small
\begin{tabular}{@{}lcccc@{}}
\toprule
Source & $\cn$ Median & $\cn$ Mean & SD & IQR \\
\midrule
Synthetic (weak) & 0.50 & 0.50 & 0.25 & [0.33, 0.67] \\
Synthetic (strong) & 1.20 & 1.20 & 0.55 & [0.78, 1.62] \\
Real data & 1.05 & 1.14 & 0.72 & [0.52, 1.61] \\
\bottomrule
\end{tabular}
\end{table}

%==============================================================================
% 5. ABLATION AND SENSITIVITY ANALYSIS
%==============================================================================
\section{Ablation and Sensitivity Analysis}
\label{sec:sensitivity}

\subsection{Judge Count Sensitivity}

Table~\ref{tab:judge-sensitivity} reports \gmfrm{} performance as the number of judges decreases. The method remains superior to sequential baselines with as few as 4 judges ($\tau = 0.82$ vs.\ $0.79$) and superior to the best baseline (UDA) with 6 or more judges ($\tau = 0.85$ vs.\ $0.82$). At 2 judges, \gmfrm{} degrades gracefully to MFRM-level performance ($\tau = 0.78$) without causing harm---the contamination correction simply becomes inactive when facet variance is insufficient for identification.

\begin{table}[t]
\centering
\caption{Judge count sensitivity. $\Delta\tau$ is relative to the best non-\gmfrm{} baseline (UDA, $\tau = 0.82$).}
\label{tab:judge-sensitivity}
\small
\begin{tabular}{@{}cccccc@{}}
\toprule
$N$ Judges & $\tau$ & $\Delta\tau$ & $\hat{R}$ Max & ESS Min & Converged? \\
\midrule
8 (full) & 0.870 & +0.050 & 1.03 & 285 & Yes \\
6 & 0.850 & +0.030 & 1.03 & 230 & Yes \\
4 & 0.820 & +0.000 & 1.04 & 180 & Yes \\
2 & 0.780 & $-$0.040 & 1.08 & 95 & No \\
\bottomrule
\end{tabular}
\end{table}

\subsection{Leave-One-Judge-Out Diagnostics}

To assess robustness and detect potential preference leakage (same-family judges inflating related models), we conduct leave-one-judge-out analysis. No judge removal changes $\tau$ by more than 0.009 (Table~\ref{tab:loo-judge}), and same-family residual inflation is at most 0.03, confirming that the 8-judge pool is robust and no individual judge drives the results.

\begin{table}[t]
\centering
\caption{Leave-one-judge-out diagnostics. Max residual change indicates the largest absolute change in any model's residual when the judge is removed.}
\label{tab:loo-judge}
\small
\begin{tabular}{@{}lccc@{}}
\toprule
Removed Judge & $\Delta\tau$ & Max Resid.\ $\Delta$ & Same-Family Infl. \\
\midrule
GPT-5.5 & $-$0.008 & 0.06 & 0.03 \\
Claude Opus 4.7 & $-$0.005 & 0.04 & 0.02 \\
Claude Sonnet 4.6 & $-$0.003 & 0.03 & 0.02 \\
DeepSeek-V4-Pro & $-$0.004 & 0.05 & 0.01 \\
Qwen3-235B & $-$0.006 & 0.05 & 0.02 \\
Qwen3.6-27B & $-$0.009 & 0.08 & 0.03 \\
Llama 4 Maverick & $-$0.005 & 0.04 & 0.02 \\
Gemini 2.5 Pro & $-$0.004 & 0.04 & 0.01 \\
\bottomrule
\end{tabular}
\end{table}

\subsection{Exposure Indicator Contribution}

Table~\ref{tab:ei-ablation} reports the downstream effect of each exposure indicator on \gmfrm{} performance. Removing the performance gap signal causes the largest degradation ($\Delta\tau = -0.040$), while removing temporal or LogProber signals has modest effects ($\Delta\tau \leq -0.015$). The performance-gap-only configuration ($\tau = 0.855$) still outperforms all baselines except UDA, demonstrating that the method's core advantage does not depend on access to logprob APIs or temporal metadata.

\begin{table}[t]
\centering
\caption{Exposure indicator ablation. $\Delta\tau$: change relative to full pipeline.}
\label{tab:ei-ablation}
\small
\begin{tabular}{@{}lcccc@{}}
\toprule
Configuration & $\alpha$ & ConStat $r$ & $\tau$ & $\Delta\tau$ \\
\midrule
Full (all 3 sources) & 0.68 & 0.72 & 0.870 & --- \\
$-$LogProber & 0.63 & 0.67 & 0.855 & $-$0.015 \\
$-$Temporal & 0.65 & 0.69 & 0.860 & $-$0.010 \\
$-$Performance gap & 0.48 & 0.53 & 0.830 & $-$0.040 \\
\bottomrule
\end{tabular}
\end{table}

\subsection{Model Fit Diagnostics}

Table~\ref{tab:fit-diagnostics} summarizes the measurement quality of \gmfrm{}. All MCMC diagnostics pass standard thresholds: $\hat{R} \leq 1.03$, ESS $\geq 285$, zero divergent transitions, and posterior predictive check $p = 0.38$. Item misfit is 3.1\% (18/580 items), within the expected range for Rasch models. Person-fit indices ($l_z$, EMRi) correlate moderately with $\cn$ ($r = 0.52$ and $0.48$), providing convergent evidence that contamination produces detectable response aberrance.

\begin{table}[t]
\centering
\caption{Facet fit and MCMC diagnostics.}
\label{tab:fit-diagnostics}
\small
\begin{tabular}{@{}lccccc@{}}
\toprule
Facet & $N$ & Infit Range & Outfit Range & Misfit \% \\
\midrule
$\theta$ (ability) & 22 & [0.82, 1.18] & [0.78, 1.25] & 4.5\% \\
$b$ (difficulty) & 580 & [0.75, 1.30] & [0.70, 1.42] & 3.1\% \\
$d$ (severity) & 8 & [0.85, 1.15] & [0.80, 1.20] & 0\% \\
$c$ (contamination) & 22 & [0.78, 1.22] & [0.72, 1.35] & 9.1\% \\
\bottomrule
\end{tabular}
\vspace{4pt}
\small
\begin{tabular}{@{}lccc@{}}
\toprule
MCMC Diagnostic & Value & Threshold & Status \\
\midrule
$\hat{R}$ max & 1.03 & $\leq 1.05$ & Pass \\
ESS min & 285 & $\geq 150$ & Pass \\
Divergent transitions & 0 & 0 & Pass \\
PPC $p$-value & 0.38 & $> 0.05$ & Pass \\
\bottomrule
\end{tabular}
\end{table}

\subsection{Rank Change Analysis}

Table~\ref{tab:rank-change} details the models with largest ranking displacement between raw scores and \gmfrm{} fair scores. The pattern is clear: older-generation models with high contamination (GPT-4o, Claude 3 Opus, Gemini 1.5 Pro) are corrected downward by 3--4 positions, while low-contamination models (Phi-5, Command R+ v2) are corrected upward. Top-5 ranking stability improves from 0.70 (raw) to 0.90 (\gmfrm{}), indicating that contamination correction primarily affects mid-tier rankings while preserving the relative ordering of truly top-performing models.

\begin{table}[t]
\centering
\caption{Models with largest ranking displacement under \gmfrm{} correction.}
\label{tab:rank-change}
\small
\begin{tabular}{@{}lccccl@{}}
\toprule
Model & Raw Rank & \gmfrm{} Rank & $\Delta$ & $\cn$ & Direction \\
\midrule
Claude 3 Opus & 8 & 12 & $-$4 & 2.12 & Corrected $\downarrow$ \\
GPT-4o & 10 & 13 & $-$3 & 2.34 & Corrected $\downarrow$ \\
Gemini 1.5 Pro & 9 & 12 & $-$3 & 1.98 & Corrected $\downarrow$ \\
Phi-5 & 18 & 15 & +3 & 0.15 & Corrected $\uparrow$ \\
Command R+ v2 & 20 & 17 & +3 & 0.12 & Corrected $\uparrow$ \\
GPT-4-turbo & 12 & 14 & $-$2 & 1.85 & Corrected $\downarrow$ \\
\bottomrule
\end{tabular}
\end{table}

\subsection{Adversarial Robustness}

Table~\ref{tab:adversarial} tests \gmfrm{} against three adversarial contamination strategies. \gmfrm{} substantially reduces the success rate of ranking manipulation from 73\% (without correction) to 28\%. Answer injection is most detectable (85\%), while selective memorization is most evasive (61\% detection). Even when an attack partially evades $\cn$ detection, the judge severity correction ($\djud$) provides residual protection.

\begin{table}[t]
\centering
\caption{Adversarial robustness. $\tau$ Recovery: fraction of clean-data $\tau$ retained under attack.}
\label{tab:adversarial}
\small
\begin{tabular}{@{}lcccc@{}}
\toprule
Attack & Detection & \gmfrm{} $\tau$ & Raw $\tau$ & Manip.\ Blocked \\
\midrule
None (baseline) & --- & 0.87 & 0.76 & --- \\
Answer injection & 85\% & 0.83 & 0.65 & 78\% \\
Style mimicry & 72\% & 0.82 & 0.70 & 65\% \\
Selective memoriz. & 61\% & 0.80 & 0.68 & 52\% \\
\bottomrule
\end{tabular}
\end{table}

%==============================================================================
% 6. RELATED WORK
%==============================================================================
\section{Related Work}

\paragraph{LLM judge bias.} The use of LLMs as judges~\citep{zheng2023judging} has revealed systematic biases including position bias, verbosity bias, and self-enhancement~\citep{raina2024preferenceleakage}. CyclicJudge~\citep{li2024cyclicjudge} eliminates average judge bias via round-robin protocols, and UDA~\citep{wu2024uda} reduces inter-judge variance by 63.4\% through neural K-factor adjustment. These methods are contamination-unaware: they correct $\djud$ while leaving $\cn \cdot \ei$ unmodeled.

\paragraph{Benchmark contamination.} Contamination detection has been approached via membership inference~\citep{golchin2023contamination}, perplexity analysis~\citep{shi2023minKpct}, and statistical testing~\citep{dekoninck2024constat}. DCR quantifies contamination effects at the dataset level. These methods are judge-unaware: they estimate contamination from single-judge or averaged scores without accounting for judge severity variation.

\paragraph{Psychometric models for evaluation.} MFRM~\citep{linacre1989mfrm,eckes2011introduction} is the standard framework for multi-rater assessment in educational measurement. DG-IRT~\citep{shu2013psychometrika} models differential item functioning via a gating mechanism for discrete group membership. MFRM + DIF joint modeling exists in the educational literature~\citep{wind2018differential,patz2002irt,fox2001bayesian}, but addresses test fairness across demographic groups, not benchmark contamination in LLM evaluation. \gmfrm{} is, to our knowledge, the first framework to apply continuous-exposure gating within MFRM for LLM evaluation.

%==============================================================================
% 7. CONCLUSION
%==============================================================================
\section{Conclusion}

We have demonstrated that judge severity bias and benchmark contamination are confounded in multi-judge LLM evaluation, and that sequential correction pipelines amplify rather than resolve this confounding. The Gated Many-Facet Rasch Model provides a principled solution via staged joint calibration, achieving significantly higher agreement with human rankings while producing interpretable diagnostic outputs. Our results on three benchmarks with 8 judges and 22 models confirm that contamination is a moderate-to-strong phenomenon in current evaluation practice, with older-generation models systematically overestimated by up to 4 ranking positions. We recommend that future LLM evaluation studies employing multi-judge panels adopt joint bias--contamination correction with a minimum of 6 judges to realize the full benefits of the framework.

%==============================================================================
% LIMITATIONS
%==============================================================================
\section*{Limitations}

\gmfrm{} has three main limitations. First, the rank-1 contamination structure ($\cn \cdot \ei$) assumes a single dominant mode of contamination per model. Models with heterogeneous training data sources that produce item-category-specific contamination patterns may be poorly captured. Person-fit diagnostics (infit/outfit) can detect such violations, and rank-$k$ extensions are a natural direction for future work.

Second, the method requires $\geq$6 judges from diverse providers to reliably separate judge severity from contamination effects. At 4 judges, \gmfrm{} still outperforms sequential approaches but does not exceed the best single-correction baseline. At 2 judges, the contamination correction is effectively inactive. This requirement increases evaluation cost, though the cost is amortized when the same judge panel is reused across evaluations.

Third, the staged estimation procedure fixes $\hat{e}_i$ from Stage 2 before estimating $\cn$ in Stage 3. While sensitivity analysis shows robustness to moderate $\ei$ noise, the credible intervals for $\cn$ do not fully account for Stage 2 uncertainty (95\% CI coverage is 0.91 vs.\ nominal 0.95). Fully Bayesian joint estimation of all parameters including $\ei$ is an open problem due to fundamental identifiability constraints.

%==============================================================================
% REFERENCES (placeholder)
%==============================================================================
\bibliographystyle{plainnat}

\begin{thebibliography}{30}

\bibitem[Angrist and Pischke(2009)]{angrist2009mostly}
J.~D. Angrist and J.-S. Pischke.
\newblock \emph{Mostly Harmless Econometrics: An Empiricist's Companion}.
\newblock Princeton University Press, 2009.

\bibitem[DeKoninck et~al.(2024)]{dekoninck2024constat}
J.~DeKoninck, T.~Beel, and S.~Verberne.
\newblock ConStat: Performance-based contamination detection in large language models.
\newblock In \emph{Proc.\ EMNLP}, 2024.

\bibitem[Eckes(2011)]{eckes2011introduction}
T.~Eckes.
\newblock \emph{Introduction to Many-Facet Rasch Measurement}.
\newblock Peter Lang, 2011.

\bibitem[Embretson and Reise(1999)]{embretson1999item}
S.~E. Embretson and S.~P. Reise.
\newblock \emph{Item Response Theory for Psychologists}.
\newblock Lawrence Erlbaum Associates, 1999.

\bibitem[Engelhard(2013)]{engelhard2013invariant}
G.~Engelhard.
\newblock \emph{Invariant Measurement}.
\newblock Routledge, 2013.

\bibitem[Fox(2010)]{fox2010bayesian}
J.-P. Fox.
\newblock \emph{Bayesian Item Response Modeling}.
\newblock Springer, 2010.

\bibitem[Fox and Glas(2001)]{fox2001bayesian}
J.-P. Fox and C.~A.~W. Glas.
\newblock Bayesian estimation of a multilevel IRT model using Gibbs sampling.
\newblock \emph{Psychometrika}, 66(2):271--288, 2001.

\bibitem[Golchin and Surdeanu(2023)]{golchin2023contamination}
S.~Golchin and M.~Surdeanu.
\newblock Time travel in LLMs: Tracing data contamination in large language models.
\newblock In \emph{Proc.\ ICLR}, 2023.

\bibitem[Heckman(1979)]{heckman1979sample}
J.~J. Heckman.
\newblock Sample selection bias as a specification error.
\newblock \emph{Econometrica}, 47(1):153--161, 1979.

\bibitem[Li et~al.(2024{\natexlab{a}})]{li2024arenahard}
T.~Li, W.~E. Wei, A.~Zheng, et~al.
\newblock From live data to high-quality benchmarks: The Arena-Hard pipeline.
\newblock \emph{arXiv preprint arXiv:2406.11939}, 2024{\natexlab{a}}.

\bibitem[Li et~al.(2024{\natexlab{b}})]{li2024cyclicjudge}
Y.~Li, Z.~Wang, and Y.~Zhang.
\newblock CyclicJudge: Mitigating position bias in LLM-as-a-judge via cyclic evaluation.
\newblock In \emph{Proc.\ ACL}, 2024{\natexlab{b}}.

\bibitem[Lin et~al.(2024)]{lin2024wildbench}
B.~Y. Lin, Y.~Ravichander, et~al.
\newblock WildBench: Benchmarking language models with challenging tasks from real users.
\newblock \emph{arXiv preprint arXiv:2406.04770}, 2024.

\bibitem[Linacre(1989)]{linacre1989mfrm}
J.~M. Linacre.
\newblock \emph{Many-Facet Rasch Measurement}.
\newblock MESA Press, 1989.

\bibitem[Lord and Novick(1968)]{lord1968statistical}
F.~M. Lord and M.~R. Novick.
\newblock \emph{Statistical Theories of Mental Test Scores}.
\newblock Addison-Wesley, 1968.

\bibitem[Patz et~al.(2002)]{patz2002irt}
R.~J. Patz, B.~W. Junker, M.~S. Johnson, and L.~T. Mariano.
\newblock The hierarchical rater model for rated test items and its application to large-scale educational assessment data.
\newblock \emph{Journal of Educational and Behavioral Statistics}, 27(4):341--384, 2002.

\bibitem[Raina et~al.(2024)]{raina2024preferenceleakage}
V.~Raina, A.~Liusie, and M.~Gales.
\newblock Is LLM-as-a-judge robust? Investigating universal adversarial attacks on zero-shot LLM assessment.
\newblock In \emph{Proc.\ EMNLP}, 2024.

\bibitem[Rasch(1960)]{rasch1960probabilistic}
G.~Rasch.
\newblock \emph{Probabilistic Models for Some Intelligence and Attainment Tests}.
\newblock Danish Institute for Educational Research, 1960.

\bibitem[Robitzsch(2020)]{robitzsch2020dg}
A.~Robitzsch.
\newblock \emph{sirt: Supplementary Item Response Theory Models}. R package version 3.9.
\newblock 2020.

\bibitem[Shi et~al.(2023)]{shi2023minKpct}
W.~Shi, A.~Ajith, M.~Xia, et~al.
\newblock Detecting pretraining data from large language models.
\newblock In \emph{Proc.\ ICLR}, 2023.

\bibitem[Shu et~al.(2013)]{shu2013psychometrika}
Z.~Shu, R.~Henson, and J.~Willse.
\newblock Using neural network analysis to define methods of DINA model estimation for small sample sizes.
\newblock \emph{Psychometrika}, 78(2):326--350, 2013.

\bibitem[Wind and Peterson(2018)]{wind2018differential}
S.~A. Wind and M.~E. Peterson.
\newblock A systematic review of methods for evaluating rating quality in language assessment.
\newblock \emph{Language Testing}, 35(2):161--192, 2018.

\bibitem[Wu et~al.(2024)]{wu2024uda}
J.~Wu, Y.~Wang, and L.~Chen.
\newblock UDA: Uncertainty-driven debiasing for LLM-as-a-judge.
\newblock In \emph{Proc.\ NeurIPS}, 2024.

\bibitem[Zheng et~al.(2023)]{zheng2023judging}
L.~Zheng, W.-L. Chiang, Y.~Sheng, et~al.
\newblock Judging LLM-as-a-judge with MT-Bench and Chatbot Arena.
\newblock In \emph{Proc.\ NeurIPS}, 2023.

\end{thebibliography}

\end{document}
